% ════════════════════════════════════════════════════════════════
\subsection{Matroidy}
% ════════════════════════════════════════════════════════════════

\begin{definition} \label{def:matroid}
	\uemph{Matroid} to para $M = (S, \mathcal{A})$ taka, że:
	\begin{enumerate}
		\item \label{matroid:skończoność} $S$ -- niepusty zbiór skończony,
		\item \label{matroid:niepustość} $\mathcal{A} \subseteq 2^S$ i $\mathcal{A} \neq \emptyset$ -- czyli $\mathcal{A}$ to niepusta rodzina podzbiorów $S$,
		\item \label{matroid:dziedziczność} $\forall A, B \subseteq S \quad A \in \mathcal{A} \wedge B \subseteq A \implies B \in \mathcal{A}$,
		\item \label{matroid:wymiana} własność wymiany:
		      \begin{flalign*}
			       & \forall A, B \subseteq S \quad A, B \in \mathcal{A} \wedge |B| > |A| \implies \text{istnieje } x \in B \setminus A \text{ taki, że } A \cup \{x\} \in \mathcal{A}. &  &  &  &
		      \end{flalign*}
	\end{enumerate}
	$\mathcal{A}$ -- rodzina zbiorów niezależnych matroidu $M$, elementy $\mathcal{A}$ -- zbiory niezależne matroidu $M$.
\end{definition}

\begin{example} \label{ex:matroid-podstawowe} \leavevmode
	\begin{enumerate}[ref=\thedefinition.\arabic*]
		\item \label{ex:matroid:wektorowy} $A$ -- macierz, $S$ -- zbiór wierszy $A$, $\mathcal{A} = \{R \subseteq S : R \text{ tworzy zbiór liniowo niezależny}\}$.\\
		      Wtedy $(S, \mathcal{A})$ to matroid.
		\item \label{ex:matroid:grafowy} $G = (V, E)$ -- graf, $\mathcal{A} = \{F \subseteq E : G[F] \text{ jest acykliczny}\}$.\\
		      Wtedy $(E, \mathcal{A})$ to matroid zwany matroidem grafowym.
	\end{enumerate}
	Z definicji $\emptyset \in \mathcal{A}$.
\end{example}

\begin{lemma} \label{lemma:matroid:liczność}
	Niech $M = (S, \mathcal{A})$. Wszystkie maksymalne zbiory niezależne $M$ mają tę samą liczność.
\end{lemma}
\begin{problem}
	$M = (S, \mathcal{A})$ -- matroid, $w: S \to \Rp$ -- funkcja wagi.\\
	Dla $A \subseteq S$ niech $w(A) = \sum_{x \in A} w(x)$ -- waga $A$.\\
	Znaleźć zbiór niezależny o maksymalnej wadze w $M$.\\
	$S = \{x_1, x_2, \ldots, x_n\}$.
\end{problem}

\begin{alg}{Algorytm Greedy}{alg:greedy}
	\FunctionNN{Greedy}{$M, w$}
	\State $n \Assign |S|$
	\State $A \Assign \emptyset$ \Comment{$A$ -- rozwiązanie}
	\State posortuj $S$ w porządku nierosnącym względem wagi, tak żeby $w(x_1) \geq w(x_2) \geq \ldots \geq w(x_n)$
	\For{$i \Assign 1$ \To $n$}
	\If{$A \cup \{x_i\} \in \mathcal{A}$} \label{line:greedy-niezaleznosc}
	\State $A \Assign A \cup \{x_i\}$
	\EndIf
	\EndFor
	\State \Return $A$
	\EndFunction
\end{alg}

\begin{theorem}[Rado, Edmonds] \label{thm:rado-edmonds}
	Algorytm Greedy znajduje zbiór niezależny o maksymalnej wadze. Co więcej, zachłanny wybór daje
	optimum dla \uemph{każdej} funkcji wagi wtedy i~tylko wtedy, gdy $(S, \mathcal{A})$ jest matroidem.
\end{theorem}

\paragraph{Złożoność czasowa algorytmu Greedy}
$n = |S|$
\begin{itemize}[nosep]
	\item linia 1: $\BO(1)$
	\item linia 2: $\BO(1)$
	\item linia 3: $\BO(n \log n)$
	\item liczba iteracji pętli z linii 4 wynosi $n$
	\item linia 5: $f(n)$ -- czas sprawdzenia warunku z linii 5
	\item linia 6: $\BO(1)$
	\item linia 7: $\BO(1)$
\end{itemize}
Pętla wykonuje $n$ iteracji, a~każda kosztuje $f(n) + \BO(1)$, co daje koszt pętli $\BO\bigl(n \cdot f(n)\bigr)$. Łącznie z~sortowaniem otrzymujemy
\[
	\BO\bigl(n \log n + n \cdot f(n)\bigr),
\]
gdzie $f(n)$ to czas sprawdzenia niezależności $A \cup \{x_i\} \in \mathcal{A}$ (linia \ref{line:greedy-niezaleznosc}). Złożoność zależy więc od reprezentacji matroidu: dla matroidów, w~których test niezależności jest wielomianowy, cały algorytm działa w~czasie wielomianowym.

% ════════════════════════════════════════════════════════════════
\subsection{Problem szeregowania zadań}
% ════════════════════════════════════════════════════════════════

\begin{itemize}[nosep]
	\item $S = \{z_1, z_2, \ldots, z_n\}$ -- zbiór zadań o jednostkowym czasie wykonania
	\item $d_1, d_2, \ldots, d_n$ -- deadliny, $d_i \in \{1, \ldots, n\}$ dla $i \in [n]$
	\item $w_1, w_2, \ldots, w_n$ -- kary, $w_i \geq 0$ dla $i \in [n]$
\end{itemize}
Płacimy karę $w_i$ za zadanie $z_i$, jeżeli nie jest skończone w momencie $d_i$.\\
\begin{problem}[szeregowania zadań z karami]
	Znaleźć kolejność wykonania zadań minimalizującą sumę kar.
\end{problem}

Załóżmy, że mamy pewien harmonogram. Zadanie $z_i$ nazywamy \uemph{terminowym} w tym harmonogramie, jeśli jest zakończone do chwili $d_i$, w przeciwnym razie nazywamy je \uemph{spóźnionym}.\\
Nie płacimy kar za zadania terminowe, płacimy kary za zadania spóźnione.

\begin{lemma}
	Każdy harmonogram można przetransformować na harmonogram z taką samą sumą kar, w którym wszystkie zadania terminowe poprzedzają wszystkie zadania spóźnione oraz zadania terminowe są posortowane w kolejności niemalejącej względem deadlinów.
\end{lemma}